% ARCHIVED DRAFT: independently audited by root and incorporated on 2026-11-19 16:37 UTC.
% The assembled paper is authoritative; this file retains drafting provenance.
\begin{saresult}{B.070}{Proved}{Every binomial base ring and every coefficient extension}
Let $R$ be a nonzero binomial ring: its additive group is
$\mathbb Z$-torsion-free, and $\binom r n\in R$ for every $r\in R$ and
$n\ge0$, with the fraction interpreted in $\mathbb Q\otimes R$.
For every unital commutative inclusion $R\subset S$, the polynomial
coaddition and comultiplication tensors are uniquely determined if and
only if the additive group of $S$ is $\mathbb Z$-torsion-free.
In every positive case all multivariable polynomial tensor maps are
isomorphisms. No Noetherian, domain or reducedness hypothesis is required
on $R$, and $S$ may have torsion as an $R$-module.

This includes $R=\operatorname{Int}(\mathbb Z)$, and also the nonreduced
non-Noetherian binomial ring $R=\mathbb Z\oplus\mathbb Q\epsilon$ with
$\epsilon^2=0$.
\end{saresult}
\begin{proof}
Write $(X)_n=X(X-1)\cdots(X-n+1)$, with $(X)_0=1$.
For every polynomial $F\in S[X]$ of degree at most $d$, and every $r\in R$,
the finite Newton identity holds in $S$:
\[
               F(r)=\sum_{n=0}^d \Delta^nF(0)\binom r n.  \tag{B70.1}
\]
This does not cancel integers in $S$. For each monomial $X^j$, the
identity $r^j=\sum_n n!\,{j\brace n}\binom r n$ holds in the
torsion-free ring $R$, where ${j\brace n}$ are the integer Stirling
coefficients. Multiplying these identities by the coefficients of $F$
and adding proves (B70.1) even if $S$ has integer torsion.
Consequently $F$ preserves $R$ if and only if its finitely many values
at $0,1,\ldots,d$ lie in $R$.

The polynomials $(X)_n$ form a monic $S$-basis, and
$\Delta^i(X)_n(0)=n!$ for $i=n$, and zero otherwise. Hence there is an
exact, split coefficient description
\[
 \Omega(R,S)=\bigoplus_{n\ge0}(R:_S n!)(X)_n.              \tag{B70.2}
\]
For several variables, the identical argument on a rectangular integer
grid gives coefficient module
$(R:_S\prod_i n_i!)$ on the basis $\prod_i(X_i)_{n_i}$.

Every localization of $R$ is again binomial. Here is the needed direct
localization argument. If a polynomial $G$ in a coefficient extension
preserves $R$, then it preserves $R[1/s]$ with values in the localized
base: induct on its degree $d$, using that
$s^dG(X)-G(sX)$ has smaller degree and still preserves $R$; evaluation
at $r/s$ gives $s^dG(r/s)\in R[1/s]$. Apply this to each
$G(X)=\binom X n$ inside $\mathbb Q\otimes R$ and allow arbitrary
localization denominators. Integer torsion-freeness is also preserved
by localization. Finally the degree-bounded preservation condition in
(B70.1) is the kernel of the finite evaluation map
\[
          S^{d+1}\longrightarrow (S/R)^{d+1}.
\]
Flatness of localization therefore proves exact localization of
$\Omega$ for every $S$, including integer torsion, and the same statement
in every number of variables. No Noetherian argument is needed.

We now analyze a local binomial ring $(A,\mathfrak n)$.
If its residue field has characteristic zero, every nonzero integer is
a unit in $A$, and ordinary integer Vandermonde determinants show that
$\Omega(A,S)=A[X]$ for every extension $S$.
If the residue characteristic is $p$, one has
$a^p-a\in pA$ for every $a\in A$: the polynomial $(X)_p$ is congruent to
$X^p-X$ modulo $p$, and $(a)_p=p!\binom a p\in pA$.
Thus every element of the local ring $A/pA$ satisfies $x^p=x$.
An element $x$ of its maximal ideal satisfies
$x(1-x^{p-1})=0$, where the second factor is a unit, so $x=0$.
The residue field is consequently $\mathbb F_p$, and
\[
                           \mathfrak n=pA.                \tag{B70.3}
\]
The element $p$ is a nonzerodivisor of $A$.
By (B70.2), all occurring coefficients lie in the ring
\[
                   C=\{s\in S:p^Ns\in A
                                  \text{ for some }N\}.  \tag{B70.4}
\]
Conversely every such element occurs, by taking $n$ with $p^N\mid n!$;
the other integer prime factors are units in $A$.
The equality $\Omega(A,S)=\Omega(A,C)$ follows.

Suppose first that $p$ is a nonzerodivisor in $S$. Then $C$ embeds into
$A[1/p]$. This overring is either $A$ or all of $A[1/p]$.
Indeed if $c=a/p^N\notin A$, then $a\notin p^NA$; choose the largest
$j<N$ for which $a\in p^jA$ and write $a=p^ju$ with $u$ a unit, using
(B70.3). Thus $c=p^{j-N}u$, so $p^{-1}\in C$.
If $C=A$, the polynomial tensor maps are the ordinary ones.
If $C=A[1/p]$, equation (B70.2) gives the free $A$-basis
$\binom X n$, and its multivariable version gives the product of these
bases. The polynomial tensor maps are isomorphisms in every arity.

For necessity of the absence of $p$-torsion, put
\[
                          P=\bigcap_{j\ge0}p^jA.
\]
The quotient $\bar A=A/P$ is a discrete valuation ring with uniformizer
$p$. To verify this without a Noetherian assumption, $p$ stays regular
modulo $P$: from $pa\in P$ cancel the regular element $p$ in each
$pa=p^{j+1}b_j$ to obtain $a\in P$. Every nonzero element of $A/P$ is a
unit times a finite power of $p$, by choosing the largest $j$ with
$a\in p^jA$. Products are therefore nonzero, and every nonzero ideal is
generated by an element having the smallest such exponent. This proves
the domain, discrete valuation and Noetherian assertions directly.

The ideal $P$ is an ideal of $C$, with $PC=P$. If $p^Nc=a\in A$ and
$x\in P$, write $x=p^{N+j}b_j$ for every $j$; then
$xc=p^jb_ja\in p^jA$ for all $j$, proving $xc\in P$.
Reduction of coefficients gives a surjection
\[
 \Omega(A,C)\longrightarrow\Omega(\bar A,C/P),\qquad
 \ker=P[X]=P\Omega(A,C).                                  \tag{B70.5}
\]
Indeed lift the finitely many coefficients of any reduced polynomial.
Its values on $A$ lie in the preimage of $A/P$, which is $A$, so the
lift is integer-valued. The kernel assertion follows from $PC=P$.

Both existence of co-operations and tensor injectivity descend through
(B70.5). For the latter, lift any reduced kernel tensor. Its polynomial
image belongs to $P[X,Y]$; subtract the tensor of ordinary monomials
with those coefficients, which reduces to zero. The remaining tensor
has zero polynomial image upstairs, hence is zero. This proves descent
without flatness of $A/P$.
If $0\ne z\in S$ satisfies $pz=0$, then $z\in C$ and its image in $C/P$
is nonzero, since $P\subset A$ and $p$ is regular in $A$. This contradicts
the necessary criterion B.063 for the discrete valuation base $\bar A$.
Thus the two co-operations cannot both exist uniquely when $S$ has
$p$-torsion.

Return to the global inclusion. If $S$ is integer torsion-free, so is
every localization $S_{\mathfrak m}$. The local positive cases above
and the proved localization of $\Omega_n$ show that every tensor map
is an isomorphism, by testing its kernel and cokernel at maximal ideals.
Conversely, any nonzero integer-torsion element has a nonzero multiple
$z$ killed by a prime $p$. Choose a maximal ideal containing its proper
$R$-annihilator. Then $p$ lies in that maximal ideal and $z$ survives
localization. Localizing the purported co-operations contradicts the
local necessity just proved.

For the examples, $\operatorname{Int}(\mathbb Z)$ is a binomial ring
because $\binom{F(X)}n$ is integer-valued whenever $F$ is.
In $\mathbb Z\oplus\mathbb Q\epsilon$, evaluating any binomial polynomial
at $a+b\epsilon$ gives its integer value at $a$ plus a rational multiple
of $\epsilon$, so this ring is binomial. It is not reduced and its
square-zero ideal $\mathbb Q\epsilon$ is not finitely generated over
the ring, proving non-Noetherianity.
\end{proof}
